% Fichier généré automatiquement par tools/generate_corrections.py.
% Source : COR/COR-03-PI/64_EPAS/64_EPAS.tex
% Statut : complet (5/5 question(s) renseignée(s), 1 reprise(s)).
\begin{corrigebox}[Corrigé extrait de la banque]
\CorrigeQuestion{1}
\par\medskip\begin{center}\captionsetup{type=figure}
 \begin{tikzpicture}[xscale=2]
\tikzset{
semilog lines/.style={thin, bleuxp}, 
semilog lines 2/.style={semilog lines,bleuxpc},
semilog half lines/.style={semilog lines 2,dotted },
semilog label x/.style={semilog lines,below,font=\tiny,black},
semilog label y/.style={semilog lines,right,font=\tiny,black}
}
\begin{scope}[yscale=1/60]
\OrdBode{20}
\semilog{-1}{2}{-80}{40}

\BodeAmp[bleuxp,thick]{-1:2}{\POAmpAsymp{1}{3.6}+\IntAmp{4}}
\BodeAmp[orangexp,ultra thick]{-1:2}{\POAmp{1}{3.6}+\IntAmp{4}}

\end{scope}
\begin{scope}[yshift=-2cm,yscale=1/150]
\UniteDegre
\OrdBode{45}
\semilog{-1}{2}{-270}{0}

\BodeArg[bleuxp,thick]{-1:2}{\POArgAsymp{4}{3.6}+\IntArg{1}}
\BodeArg[orangexp,ultra thick]{-1:2}{\POArg{4}{3.6}+\IntArg{1}}
\end{scope}
\end{tikzpicture}
\end{center}\par\medskip
\CorrigeQuestion{2}
Quelle que soit la valeur du retard, il n'a aucune influence sur le gain. La phase du retard est de $-0,2\omega$ soit $-\SI{0,2}{rad} = -11\degres$ en $\omega = \SI{1}{rad/s}$ et $-\SI{2}{rad} = -114\degres$ en $\omega = \SI{10}{rad/s}$.

\marginnote{Mauvais tracé...}
\par\medskip\begin{center}\captionsetup{type=figure}
 \begin{tikzpicture}[xscale=2]
\tikzset{
semilog lines/.style={thin, bleuxp}, 
semilog lines 2/.style={semilog lines,bleuxpc},
semilog half lines/.style={semilog lines 2,dotted },
semilog label x/.style={semilog lines,below,font=\tiny,black},
semilog label y/.style={semilog lines,right,font=\tiny,black}
}
\begin{scope}[yscale=1/60]
\OrdBode{20}
\semilog{-1}{2}{-20}{20}

\BodeAmp[bleuxp,thick]{-1:2}{\RetAmp{.2}}
\BodeAmp[orangexp,ultra thick]{-1:2}{\RetAmp{.2}}

\end{scope}
\begin{scope}[yshift=-1cm,yscale=1/150]
\UniteDegre
\OrdBode{30}
\semilog{-1}{2}{-90}{0}

\BodeArg[orangexp,ultra thick]{-1:2}{\RetArg{0.2}}
\end{scope}
\end{tikzpicture}
\end{center}\par\medskip
\CorrigeQuestion{3}
Lorsque la phase vaut $-130\degres$, la gain vaut $-\SI{4}{dB}$. Il faut donc augmenter le gain de 4 dB soit $K_C = 10^{4/20} \simeq 1,6$. 

\par\medskip\begin{center}\captionsetup{type=figure}
 \begin{tikzpicture}[xscale=2]
\tikzset{
semilog lines/.style={thin, bleuxp}, 
semilog lines 2/.style={semilog lines,bleuxpc},
semilog half lines/.style={semilog lines 2,dotted },
semilog label x/.style={semilog lines,below,font=\tiny,black},
semilog label y/.style={semilog lines,right,font=\tiny,black}
}
\begin{scope}[yscale=1/60]
\OrdBode{20}
\semilog{-1}{2}{-80}{40}

\BodeAmp[bleuxp,thick]{-1:2}{\POAmpAsymp{1}{3.6}+\IntAmp{4}}
\BodeAmp[orangexp,ultra thick]{-1:2}{\POAmp{1}{3.6}+\IntAmp{4}}

\end{scope}
\begin{scope}[yshift=-2cm,yscale=1/150]
\UniteDegre
\OrdBode{45}
\semilog{-1}{2}{-270}{0}

\BodeArg[bleuxp,thick]{-1:2}{\POArgAsymp{4}{3.6}+\IntArg{1}}
\BodeArg[orangexp,ultra thick]{-1:2}{\POArg{4}{3.6}+\IntArg{1}+\RetArg{.2}}
\end{scope}
\end{tikzpicture}
\end{center}\par\medskip
\CorrigeQuestion{4}
Le gain est nul pour une pulsation de $\SI{1,5}{rad/s}$ On veut que $T_c$ n'affecte pas la marge de gain. On place donc $T_c$ une décade plus tôt, c'est à dire tel que $1/T_c = 1,5/10$ soit $T_c = \SI{6,7}{s}.$
\CorrigeQuestion{5}
\CorrigeSource{Réponse reprise d'un exercice homonyme de la banque ; la provenance exacte est indiquée dans le commentaire du fichier.}
% Réponse reprise de : PERF/PERF-02-Marges/64_EPAS/64_EPAS.tex
La BO est de classe 1, l'erreur de trainage est $\varepsilon_t = {T_c}{4K_c}$.
\end{corrigebox}
